\section{Universal Refs}

\begin{frame}[fragile]{Bounded Vector}
  \begin{sqt:result}
    When a caller passes an rvalue into a function,
    by convention, this means that the caller is giving permission 
    for the callee to \emph{move construct} from the ref.
  \end{sqt:result}
  In this code example we define a vector class with a fixed capacity.
  \begin{minted}{cpp}
    /* 3.01 */
    template <typename T, std::size_t Capacity>
    struct bounded_vector {
        // lifetime and resource management
        bounded_vector() {/* */}
        ~bounded_vector() {/* */}
        bounded_vector(const bounded_vector<T, Capacity>& other) {/* */}
        // member functions
        void push_back(T& obj) {/* */}
        void push_back(T&& obj) {/* */}
    };
  \end{minted}
\end{frame}

\begin{frame}[fragile]{Why do we want the rvalue overload?}
  What would happen if we didn't have this rvalue overload for \code{push\_back}?

  \begin{minted}{cpp}
    int main() {
        bounded_vector<intbox, 8> vec;
        vec.push_back(intbox(3)); // COMPILE ERROR

        // required either one of these to be defined:
        // - push_back(const intbox& obj) [because const lvalue refs can bind to rvalues]
        // - push_back(intbox&& obj)
    }
  \end{minted}
\end{frame}

\begin{frame}[fragile]{Why do we want the rvalue overload?}
  What would happen if we made \code{push\_back} take a \emph{const} ref instead?
  \begin{minted}{cpp}
    /* 3.01 */
    template <typename T, std::size_t Capacity>
    struct bounded_vector {
        /* ... */
        void push_back(const T& obj) {
            // this function is NOT permitted to move from obj
            // since it is a const reference
        }
    };

    int main() {
        bounded_vector<intbox, 8> vec;
        vec.push_back(intbox(3));
        // copy constructs into the vector
        // this may be expensive if type T handled more resources
    }
  \end{minted}
\end{frame}

\begin{frame}[fragile]{The issue: repeated logic}
  For this particular use case (\code{push\_back()}), we run the same logic
  when both rvalues and lvalues are passed as arguments. 
  We just want to \emph{forward the argument perfectly} into the constructor.

  \begin{minted}{cpp}
    /* 3.01 - bounded_vector */
    void push_back(T& obj) {
        assert(size_ < Capacity);
        new (data_ + size_) T(obj);             // calls constructor with obj (lvalue)
        size_++;
    }

    void push_back(T&& obj) {
        assert(size_ < Capacity);
        new (data_ + size_) T(std::move(obj));  // calls constructor with obj (rvalue)
        size_++;
    }

    // note that even though decltype of obj is T&&
    // the expression "obj" is still lvalue category, so we must std::move to make rvalue
  \end{minted}
\end{frame}

\begin{frame}[fragile]{Invalid solution: attempt to deduce \code{F = T\&}}
  \begin{minted}{cpp}
    template <typename F>                             // e.g. hope that: F = T&&
    void push_back(F obj) {                           // e.g. hope that: T&& obj
        assert(size_ < Capacity);
        new (data_ + size_) T(static_cast<F>(obj));   // e.g. static_cast<T&&>(obj)
        size_++;
    }
    /* ... */
    vec.push_back(intbox(3));
  \end{minted}

  The idea is that we hope \code{F} can deduce to either \code{T\&} or \code{T\&\&},
  depending on what value category we use at the call-site.

  \begin{sqt:result}
    Template argument deduction will strip references if the argument is declared as \code{(F obj)}.
    i.e. \code{F} will always deduce to just \code{T}, not \code{T\&} or \code{T\&\&}.
  \end{sqt:result}
\end{frame}

\begin{frame}[fragile]{Bad solution: explicit instantiation}
  \begin{minted}{cpp}
    template <typename F>                             // e.g. F = T&&
    void push_back(F obj) {                           // e.g. T&& obj
        assert(size_ < Capacity);
        new (data_ + size_) T(static_cast<F>(obj));   // e.g. static_cast<T&&>(obj)
        size_++;
    }
    /* ... */
    vec.push_back<intbox &&>(intbox(3));
  \end{minted}

  We have gotten rid of duplicate logic, at the cost of needing each call-site
  to explicitly specify the template parameter. If you miss one, it silently
  makes a by-value overload (which is the worst thing ever).  
\end{frame}

\begin{frame}[fragile]{The solution: forwarding references}
  \begin{sqt:definition}
    The \emph{forwarding reference} technique is when a function parameter has declared type
    \code{F\&\&}, where \code{F} is a template argument of the function.
  \end{sqt:definition}

  \begin{minted}{cpp}
    /* 3.02 - bounded_vector */
    template <typename F>
    void push_back(F&& obj) {
        assert(size_ < Capacity);
        new (data_ + size_) T(static_cast<F&&>(obj));
        size_++;
    }
  \end{minted}

  At first glance, wouldn't this only produce an rvalue overload?
\end{frame}

\begin{frame}[fragile]{How does it work}
  There is no special compiler logic that handles forwarding references.\\
  Two parts of the C++ language come together to make forwarding references possible:
  \begin{itemize}
    \item Reference collapsing rules
    \item A special rule for template argument deduction
  \end{itemize}

  Ref collapsing rules determine the resultant type when you take a ref to a ref.

  \begin{minted}{cpp}
      using L = int&;     // int &      -> int&
        using LL = L&;    // int & &    -> int&
        using LR = L&&;   // int & &&   -> int&   (relevant to forward refs)

      using R = int&&;    // int &&     -> int&&  (relevant to forward refs)
        using RL = R&;    // int && &   -> int&
        using RR = R&&;   // int && &&  -> int&&
  \end{minted}
\end{frame}


\begin{frame}[fragile]{How does it work}
  When called with both value categories, two instantiations are created:
  \begin{minted}{cpp}
    template <> // F = T
    void push_back(T && obj) {
        assert(size_ < Capacity);
        new (data_ + size_) T(static_cast<T &&>(obj));
        size_++;
    }

    template <> // F = T&
    void push_back(T& && obj) {                          // F&&  -> T & && -> T&
        assert(size_ < Capacity);
        new (data_ + size_) T(static_cast<T& &&>(obj));  // F&&  -> T & && -> T&
        size_++;
    }
  \end{minted}

  \begin{sqt:definition}
    \code{std::forward<F>} is defined as \code{static\_cast<F\&\&>}.
  \end{sqt:definition}
\end{frame}

